% Lösungen Einheit 3 von 4 – Multiplizieren und Dividieren (zusammenbau v0.1)
\einheitenkopf{Einheit 3 von 4 · Multiplizieren und Dividieren}

\erg{20}{a) minus \quad b) $-35$ \quad c) $42$ \quad d) $-7$ \quad e) $30$ \quad f) $(-4)^2 = 16$; $-4^2 = -16$ \quad g) $-7$ \quad h) $(15 + (-6)) : (-4) = 9 : (-4) = -2{,}25$ \quad i) $(3 + (-9)) : (-2) = (-6) : (-2) = 3$ \quad j) $3 \cdot (-4 - 5) = 3 \cdot (-9) = -27$ \quad k) immer positiv}
\erg{21}{a) nein: drei Minuszeichen sind eine ungerade Anzahl, das Produkt ist negativ}
\erg{22}{a) Sie hat das Vorzeichen vergessen; verschiedene Vorzeichen ergeben minus: $(-6) \cdot 5 = -30$. \quad b) Das Ergebnis wird jedes Mal um $4$ größer; also $(-1) \cdot (-4) = 4$, positiv. \quad c) $6 \cdot (-15) = -90$ €; der Kontostand sinkt um $90$ €.}
